\documentclass[10pt,conference]{IEEEtran}

\usepackage[utf8]{inputenc}
\usepackage[T1]{fontenc}
\usepackage[english]{babel}
\usepackage{amsmath,amssymb}
\usepackage{graphicx}
\usepackage{booktabs}
\usepackage{multirow}
\usepackage{cite}
\usepackage{url}
\usepackage{hyperref}
\usepackage{tabularx}
\usepackage{siunitx}

\hypersetup{
    colorlinks=true,
    linkcolor=blue,
    citecolor=blue,
    urlcolor=blue,
    pdftitle={Objective and Comparative Characterization of Magnetic Tape Emulations},
}

\begin{document}

\title{Objective and Comparative Characterization of Magnetic Tape Emulations: \\
Calibration Protocols, Response Maintenance vs. Gain, and Nonlinear Frequency Compression}

\author{
\IEEEauthorblockN{Alioth Ponce Aranda}
\IEEEauthorblockA{APA Audio Labs\\
Research in Acoustics and Digital Signal Processing (DSP)\\
Email: \texttt{aliothponce@gmail.com}}
}

\maketitle

\begin{abstract}
An empirical and objectively reproducible measurement protocol is established to evaluate and
compare free magnetic tape emulation plugins. The \emph{Iron Deck} plugin (M~Media Audio)
is analyzed against a representative reference processor from an internationally established
analog recording hardware firm.
The research examines the divergence of equalization standards between both devices—coinciding
only in NAB, while Iron Deck offers IEC/AES and the reference exposes CCIR—and evaluates physical
calibration protocols (MRL/STL calibration tapes, magnetic fluxes of $200$, $250$, $355$, and
$510\,\text{nWb/m}$) and their correspondence in the digital domain.
Through Python scripts, three exclusive analyses are performed: (1)~spectral calibration
stability under gain excursion (\emph{drive sweep}); (2)~static transfer curve
Input$\to$Output at multiple frequencies; and (3)~frequency-dependent nonlinear compression.
The results show that Iron Deck matches the digital integrity of the reference device
in frequency stability and anti-aliasing, and offers versatility advantages in the selection
of equalization standard and transport speed.
This study sets a methodological precedent for the audio engineering community.
\end{abstract}

\begin{IEEEkeywords}
magnetic tape emulation, MRL calibration tape, NAB/IEC/CCIR equalization, frequency response,
total harmonic distortion (THD), frequency-dependent nonlinear compression, aliasing,
audio DSP, objective analysis.
\end{IEEEkeywords}

% ===========================================================================
\section{Introduction}
% ===========================================================================

In the field of audio signal processing, analog emulations play a critical role by imparting
dynamic cohesion and harmonic coloration characteristic of the magnetic medium. Historically,
the appreciation of these processors has been subject to the commercial prestige of the developing
brand or the name of the modeled physical machine, relegating objective analysis to the background.

This research addresses a rigorous comparison between two free magnetic tape emulations from
opposite origins: \emph{Iron Deck} (M~Media Audio), a plugin from an independent developer with
parametric control of speed and equalization curve; and a reference plugin distributed by a leading
firm in analog recording hardware, treated anonymously here for reasons of technical independence.

The central objective is not to subjectively determine which device produces a more appreciated
sound coloration, but to build an empirical precedent based on reproducible DSP metrics:
frequency and phase response, total harmonic distortion (THD), frequency-dependent nonlinear
compression, spectral calibration drift, and aliasing immunity. All measurements were performed
programmatically to guarantee their reproducibility independent of perceptual criteria.

% ===========================================================================
\section{Physical Foundations and Calibration Standards}
% ===========================================================================

\subsection{Magnetic Losses and Head Geometry}

The physics of tape transport imposes restrictions dependent on speed $v$ and
magnetic wavelength $\lambda = v/f$. At low speeds ($7.5\,\text{ips}$),
the wavelength of the highs ($10\,\text{kHz}$) is drastically reduced, increasing
gap losses:
\begin{equation}
L_g(f) = 20\log_{10}\!\left|\frac{\sin(\pi g/\lambda)}{\pi g/\lambda}\right| \quad (\text{dB})
\end{equation}
where $g$ is the physical width of the playback head gap. The Wallace separation loss
\cite{wallace}, $L_d(f)=54.6\,(d/\lambda)\,\text{dB}$, attenuates the highs when the tape
separates from the head. In the low range, the geometry of the magnetic poles introduces
constructive/destructive interference known as \emph{head bump}, characteristic in the
$50$–$100\,\text{Hz}$ range.

\subsection{Equalization Standards: Convergence in NAB and Divergence in IEC/CCIR}

A fundamental finding of this study is the divergence of equalization standards between
the two evaluated devices:

\begin{itemize}
    \item \textbf{NAB} (National Association of Broadcasters): the only standard present in
    \emph{both} plugins. It uses time constants $\tau_1=50\,\mu\text{s}$ ($f_1\approx 3183\,\text{Hz}$)
    and $\tau_2=3180\,\mu\text{s}$ ($f_2\approx 50\,\text{Hz}$), with a high-frequency boost and low-frequency
    emphasis (American standard).

    \item \textbf{IEC~1\,/\,CCIR} (European standard, $15\,\text{ips}$): available only in the
    reference device. Constant $\tau_1=35\,\mu\text{s}$ ($f_1\approx 4547\,\text{Hz}$), without low-frequency
    emphasis, producing a flatter and more transparent profile in the mid-high band.

    \item \textbf{IEC~2\,/\,AES} ($30\,\text{ips}$): available only in Iron Deck.
    Constant $\tau_1=17.5\,\mu\text{s}$ ($f_1\approx 9095\,\text{Hz}$), designed for
    maximum extended response in high frequencies and lower noise floor at high speed.
\end{itemize}

The inclusion of CCIR in the reference reflects the modeling of European studio equipment
(e.g., Studer, Telefunken), while the IEC/AES offering in Iron Deck targets high-resolution
workflows at $30\,\text{ips}$. This design divergence is in itself a subject of technical study.

\subsection{Physical Calibration Protocols and their Digital Mapping}

In physical tape machines, studio alignment requires certified calibration tapes (MRL/STL) \cite{mrl}
to adjust:
\begin{enumerate}
    \item \textbf{Azimuth}: perpendicularity of the playback head using a $10$–$15\,\text{kHz}$ tone.

    \item \textbf{Operating flux}: calibration of $0\,\text{VU}\equiv+4\,\text{dBu}\equiv -18\,\text{dBFS}$
    with a $1\,\text{kHz}$ tone at fluxes of $200$, $250$, $355$, or $510\,\text{nWb/m}$.

    \item \textbf{Bias adjustment}: polarization current of $100$–$200\,\text{kHz}$ adjusted by \emph{overbias}
    ($-1.5\,\text{dB}$ at $10\,\text{kHz}$) to optimize the balance between distortion and noise.
\end{enumerate}

In the digital domain, this protocol is evaluated by verifying that the processing does not introduce
frequency drifts or erratic phase shifts under gain excursions, which is the subject of
Section~\ref{sec:calibracion} of this work.

% ===========================================================================
\section{Programmatic DSP Analysis Methodology}
% ===========================================================================

Measurements were performed programmatically in Python~3 (\texttt{scipy.signal}, \texttt{numpy})
at $48\,\text{kHz}$ / 32-bit float, on signals generated and processed by each plugin in real-time
under the exact conditions of each experiment:

\begin{itemize}
    \item \textbf{Frequency response}: logarithmic sine sweep ($1\,\text{Hz}$–$24\,\text{kHz}$).
    The transfer function $H(f)$ was estimated using the cross-spectral density (CSD) ratio:
    \begin{equation}
        H(f) = \frac{S_{xy}(f)}{S_{xx}(f)}
    \end{equation}
    with Savitzky-Golay smoothing ($w=151$, $p=2$) to remove spectral noise without distorting the EQ slopes.

    \item \textbf{Static transfer and THD}: pure tones at $100\,\text{Hz}$, $1\,\text{kHz}$, and $5\,\text{kHz}$,
    evaluated at input levels of $-30$, $-18$, $-12$, $-6$, and $0\,\text{dBFS}$ to characterize compression
    and total harmonic distortion ($H_2$–$H_{19}$).

    \item \textbf{Calibration stability}: reference tones at $3000\,\text{Hz}$ (NAB pattern tone) and
    $3150\,\text{Hz}$ (DIN/IEC), measured with Welch ($N_{\text{FFT}}=4f_s$) for sub-Hz resolution.

    \item \textbf{Spectral distribution}: pink noise processed and segmented into 10 octave bands
    centered from $31.5$–$16000\,\text{Hz}$.
\end{itemize}

Input files follow the naming convention:
\begin{center}
\texttt{\{Device\}-\{IPS\}\_IPS-EQ\_\{curve\}-\{P1\}\_\{V1\}-\{P2\}\_\{V2\}.wav}
\end{center}
which allows automatic extraction of metadata (speed, EQ standard, drive level) via regular expressions,
without assuming a fixed parameter grid.

% ===========================================================================
\section{Experimental Results and Analysis}
% ===========================================================================

\subsection{Nominal Tonal Portrait: Iron Deck}
\label{sec:retrato_iron}

Fig.\,\ref{fig:01a_retrato_iron} presents the direct comparison of the spectral response $H(f)$ at
nominal level ($-12\,\text{dBFS}$) of Iron Deck between NAB ($15\,\text{ips}$) and IEC ($30\,\text{ips}$) modes.
The NAB curve exhibits a \emph{head bump} of approximately $+0.8\,\text{dB}$ at $75\,\text{Hz}$,
characteristic of the NAB bass boost, with a progressive roll-off above $10\,\text{kHz}$.
The IEC curve at $30\,\text{ips}$ extends the linear response up to $22\,\text{kHz}$, with a considerably
flatter behavior in the mid-high band. The differential panel $\Delta H(f)$ quantitatively isolates
the gain or loss per band when switching between standards.

\begin{figure}[htbp]
\centering
\includegraphics[width=0.48\textwidth]{resultados_analisis/01a_Retrato_Tonal_IronDeck.png}
\caption{Nominal tonal portrait of Iron Deck. Comparison between NAB curve ($15\,\text{ips}$)
and IEC ($30\,\text{ips}$) at $-12\,\text{dBFS}$ input level. The lower panel shows
the spectral difference $\Delta H(f)$ (IEC$-$NAB) in dB.}
\label{fig:01a_retrato_iron}
\end{figure}

\subsection{Nominal Tonal Portrait: Reference Device}
\label{sec:retrato_ref}

Fig.\,\ref{fig:01b_retrato_ref} presents the tonal portrait of the reference device comparing NAB
($15\,\text{ips}$) and CCIR ($15\,\text{ips}$). The CCIR curve shows progressive attenuation above
$12\,\text{kHz}$ and the absence of a bass boost, characteristic features of the European standard.
The difference $\Delta H(f)$ relative to NAB exceeds $-3\,\text{dB}$ in the high-frequency bands,
quantifying the tonal impact of the playback standard choice.

\begin{figure}[htbp]
\centering
\includegraphics[width=0.48\textwidth]{resultados_analisis/01b_Retrato_Tonal_Referencia.png}
\caption{Nominal tonal portrait of the reference device. Comparison between NAB and
CCIR (both at $15\,\text{ips}$, $-12\,\text{dBFS}$). The lower panel shows
$\Delta H(f)$ (CCIR$-$NAB) in dB.}
\label{fig:01b_retrato_ref}
\end{figure}

Figs.\,\ref{fig:02_velocidad} and \ref{fig:03_eq} detail the effect of transport speed ($7.5$, $15$,
$30\,\text{ips}$) and equalization curve (NAB vs.\ IEC) in Iron Deck, respectively.

\begin{figure}[htbp]
\centering
\includegraphics[width=0.48\textwidth]{resultados_analisis/02_Efecto_Velocidad_Cinta_IronDeck.png}
\caption{Effect of transport speed in Iron Deck ($7.5$, $15$, and $30\,\text{ips}$)
under NAB equalization, at $-12\,\text{dBFS}$ input. At lower speeds, high-frequency gap losses
and the low-frequency head bump are accentuated.}
\label{fig:02_velocidad}
\end{figure}

\begin{figure}[htbp]
\centering
\includegraphics[width=0.48\textwidth]{resultados_analisis/03_Efecto_Curva_EQ_IronDeck.png}
\caption{Comparison between NAB and IEC equalizations in Iron Deck at $30\,\text{ips}$.
The IEC/AES standard produces a more extended response in high frequencies and lacks
the bass emphasis typical of NAB.}
\label{fig:03_eq}
\end{figure}

\subsection{Exclusive Analysis: Calibration Maintenance vs. Gain}
\label{sec:calibracion}

To determine whether the saturation (\emph{drive}) control alters the plugin's frequency calibration,
the relative tonal drift is defined:
\begin{equation}
\Delta H(f,\,D) = |H(f,D)|_{\text{dB}} - |H(f,D_{\text{nom}})|_{\text{dB}}
\end{equation}
where $D_{\text{nom}}$ is the nominal operating point (REC control~=~5 for Iron Deck,
Input~=~0~dB for the reference). Fig.\,\ref{fig:09_calibracion} presents this drift for both
devices at $15\,\text{ips}$, NAB.

\begin{figure}[htbp]
\centering
\includegraphics[width=0.48\textwidth]{resultados_analisis/09_Mantenimiento_Calibracion_vs_Ganancia.png}
\caption{Maintenance of spectral calibration as a function of the drive level
(upper panel: $H(f)$ per drive point; lower panel: drift $\Delta H(f)$ relative to
the nominal point). Left: Iron Deck (REC 0 / 5 / 10). Right: Reference
(Input $-12$ / $0$ / $+24$~dB).}
\label{fig:09_calibracion}
\end{figure}

The results confirm that Iron Deck maintains the spectral response calibration within a margin
of $|\Delta H(f)| < 0.15\,\text{dB}$ for any drive point in the $20\,\text{Hz}$–$20\,\text{kHz}$ band,
adding harmonic saturation without altering the spectral balance. The reference device exhibits
a drift greater than $+24\,\text{dB}$ of input, with high-frequency compression of up to
$\sim$$0.5\,\text{dB}$, resulting from the limit of the magnetization curve at extreme excitation.

\subsection{Exclusive Analysis: Static Transfer Curve}
\label{sec:transferencia}

Table\,\ref{tab:transferencia} synthesizes the \emph{relative level compression}
($\Delta_{\text{REC}} = \overline{\Delta_C}(\text{REC}_{\text{max}}) -
\overline{\Delta_C}(\text{REC}_{\text{min}})$) measured at $1\,\text{kHz}$ and $100\,\text{Hz}$
for each evaluated drive point. $\Delta_C = \text{Out}_{\text{RMS}} -
\text{In}_{\text{RMS}}$ (in dB) quantifies the deviation of the insertion gain relative
to the linear 1:1 relationship; $\Delta_{\text{REC}}$ is the total compression excursion
when sweeping the full range of the drive control.

\begin{table}[htbp]
\caption{Compression excursion when traversing the entire drive range.
$\Delta_C = \text{Out}_\text{RMS} - \text{In}_\text{RMS}$ (dB) measured at
$-3\,\text{dBFS}$ input. $\Delta_\text{REC}$ is the difference in $\Delta_C$ between the
maximum and minimum evaluated drives ($15\,\text{ips}$, NAB).}
\label{tab:transferencia}
\centering
\small
\begin{tabular}{llcrr}
\toprule
Device & Drive & $f$ & $\Delta_C$ (dB) & $\Delta_\text{REC}$ (dB) \\
\midrule
Iron Deck & REC 0  & 1\,kHz & ref.\,$\Delta_C^0$      & \multirow{3}{*}{$-19.0$} \\
Iron Deck & REC 5  & 1\,kHz & $\Delta_C^0 - 3.67$  & \\
Iron Deck & REC 10 & 1\,kHz & $\Delta_C^0 - 18.99$ & \\
\midrule
Iron Deck & REC 0  & 100\,Hz & ref.\,$\Delta_C^0$      & \multirow{3}{*}{$-19.1$} \\
Iron Deck & REC 5  & 100\,Hz & $\Delta_C^0 - 3.74$  & \\
Iron Deck & REC 10 & 100\,Hz & $\Delta_C^0 - 19.08$ & \\
\midrule
Reference & Input $-12$\,dB & 1\,kHz  & ref.\,$\Delta_C^0$      & \multirow{3}{*}{$-21.8$} \\
Reference & Input $+0$\,dB  & 1\,kHz  & $\Delta_C^0 - 0.03$  & \\
Reference & Input $+24$\,dB & 1\,kHz  & $\Delta_C^0 - 21.85$ & \\
\midrule
Reference & Input $-12$\,dB & 100\,Hz & ref.\,$\Delta_C^0$      & \multirow{3}{*}{$-22.7$} \\
Reference & Input $+0$\,dB  & 100\,Hz & $\Delta_C^0 - 0.47$  & \\
Reference & Input $+24$\,dB & 100\,Hz & $\Delta_C^0 - 22.66$ & \\
\bottomrule
\end{tabular}
\end{table}

Table\,\ref{tab:transferencia} shows that the compression excursion over the entire drive range
is $\approx 19\,\text{dB}$ for Iron Deck and $\approx 21$--$23\,\text{dB}$ for the reference,
with similar behavior at $1\,\text{kHz}$ and $100\,\text{Hz}$.

\subsection{Exclusive Analysis: Nonlinear Compression by Frequency}
\label{sec:compresion_frec}

Table\,\ref{tab:compresion} presents the analysis of the variation of $\Delta_C$ when sweeping the
input level range ($-33$ to $-3\,\text{dBFS}$) at the nominal drive setting, for three frequency bands.
The ``Excursion'' column reports $\delta = \Delta_C^{\text{max}} - \Delta_C^{\text{min}}$,
indicating by how many dB the effective compression varies when sweeping the entire dynamic range:
a greater excursion implies a more nonlinear transfer curve.

\begin{table}[htbp]
\caption{Variation of net compression $\Delta_C = \text{Out}_\text{RMS} -
\text{In}_\text{RMS}$ when sweeping the input range from $-33$ to $-3\,\text{dBFS}$,
at nominal drive setting ($15\,\text{ips}$, NAB). The excursion
$\delta = \Delta_C^{\text{max}} - \Delta_C^{\text{min}}$ indicates the degree
of nonlinearity of the transfer curve in each band.}
\label{tab:compresion}
\centering
\small
\begin{tabular}{lrrrr}
\toprule
Device & $f$ (Hz) & $\Delta_C^{\text{min}}$ (dB) & $\Delta_C^{\text{max}}$ (dB) & $\delta$ (dB) \\
\midrule
Iron Deck  & 100  & $-3.25$ & $-0.02$ & $3.23$ \\
Iron Deck  & 1000 & $-3.26$ & $-0.00$ & $3.26$ \\
Iron Deck  & 5000 & $-3.25$ & $+0.01$ & $3.26$ \\
\midrule
Reference & 100  & $-1.93$ & $-0.00$ & $1.93$ \\
Reference & 1000 & $-1.44$ & $-0.00$ & $1.44$ \\
Reference & 5000 & $-1.14$ & $+0.02$ & $1.16$ \\
\bottomrule
\end{tabular}
\end{table}

The data in Table\,\ref{tab:compresion} reveal that Iron Deck exhibits a greater compression
excursion ($\delta \approx 3.25\,\text{dB}$) that is uniform across the three bands, consistent
with a symmetric \emph{soft-knee} magnetization curve. In the reference, nonlinearity decreases
with frequency ($\delta = 1.93\,\text{dB}$ at $100\,\text{Hz}$, $1.16\,\text{dB}$ at $5\,\text{kHz}$),
documenting the physical effect of high-frequency recording pre-emphasis that reduces the flux
excursion and, therefore, saturation.

\subsection{THD, Anti-Aliasing, and Spectral Distribution}
\label{sec:thd}

Fig.\,\ref{fig:05_thd} plots THD (\%) at $1\,\text{kHz}$ as a function of input level,
dominated by odd harmonics ($H_3$, $H_5$) in Iron Deck, consistent with the Preisach-Jiles-Atherton
magnetic hysteresis model. Fig.\,\ref{fig:06_aliasing} demonstrates the effectiveness of
Iron Deck's oversampling filter: at $5\,\text{kHz}$ and $0\,\text{dBFS}$, the harmonic peaks
are found exactly at multiples of the fundamental without spurious energy between them,
evidencing the absence of \emph{foldover} up to the Nyquist frequency.

\begin{figure}[htbp]
\centering
\includegraphics[width=0.48\textwidth]{resultados_analisis/05_THD_vs_Nivel.png}
\caption{Total Harmonic Distortion (THD\,\%) vs.\ input level at $1\,\text{kHz}$,
$15\,\text{ips}$, NAB. Left: Iron Deck (three REC drive points). Right:
reference device (three Input levels).}
\label{fig:05_thd}
\end{figure}

\begin{figure}[htbp]
\centering
\includegraphics[width=0.48\textwidth]{resultados_analisis/06_Aliasing_5kHz_Saturacion_Extrema.png}
\caption{Anti-aliasing test with a $5\,\text{kHz}$ tone at maximum saturation
($0\,\text{dBFS}$). Magnitude normalized to the fundamental. Vertical markers indicate
multiples of $5\,\text{kHz}$. The absence of spurious energy between harmonics confirms
effective oversampling filtering.}
\label{fig:06_aliasing}
\end{figure}

Fig.\,\ref{fig:07_psd} presents the energy distribution in octave bands under pink noise
excitation, with its own scaled differential panel to detect differences of tenths of a dB
that would go unnoticed in superimposed curves.

\begin{figure}[htbp]
\centering
\includegraphics[width=0.48\textwidth]{resultados_analisis/07_Distribucion_Espectral_RuidoRosa.png}
\caption{Spectral energy distribution in octave bands (pink noise, $-12\,\text{dBFS}$,
$15\,\text{ips}$, NAB). Left panel: Iron Deck vs.\ Reference at nominal setting. Right panel:
Iron Deck at minimum vs.\ maximum drive (spectral migration). The lower differential panel
in each case isolates the band-by-band variation in dB.}
\label{fig:07_psd}
\end{figure}

\subsection{Digital Frequency Stability}
\label{sec:estabilidad}

Table\,\ref{tab:estabilidad} summarizes the digital stability test using pattern tones of
$3000\,\text{Hz}$ (NAB) and $3150\,\text{Hz}$ (DIN/IEC), processed at all available drive points.
The detected frequency was measured with Welch ($N_{\text{FFT}} = 4f_s$, resolution $0.25\,\text{Hz}$).

\begin{table}[htbp]
\caption{Digital Frequency Stability Test. The detected frequency was measured
with $0.25\,\text{Hz}$ resolution (Welch, $N_\text{FFT}=4f_s$). A representative row
is reported per device and standard combination; all evaluated drive points
yielded $\text{Drift}=+0.00\,\text{Hz}$.}
\label{tab:estabilidad}
\centering
\small
\begin{tabular}{llccc}
\toprule
Device & EQ & IPS & Pattern (Hz) & Drift (Hz) \\
\midrule
Iron Deck & NAB & 7.5 & 3000 & $+0.00$ \\
Iron Deck & NAB & 15  & 3000 & $+0.00$ \\
Iron Deck & NAB & 30  & 3000 & $+0.00$ \\
Iron Deck & IEC & 30  & 3150 & $+0.00$ \\
\midrule
Reference & NAB  & 7.5 & 3000 & $+0.00$ \\
Reference & NAB  & 15  & 3000 & $+0.00$ \\
Reference & CCIR & 15  & 3150 & $+0.00$ \\
\bottomrule
\end{tabular}
\end{table}

Both devices demonstrate zero frequency drift ($+0.00\,\text{Hz}$) under all evaluated drive points,
confirming that neither Iron Deck nor the reference device introduces frequency modulation or
digital clock artifacts.

% ===========================================================================
\section{Discussion}
% ===========================================================================

\begin{enumerate}
    \item \textbf{Spectral response at $30\,\text{ips}$ (IEC/AES):} Iron Deck offers a linear
    extension up to $22\,\text{kHz}$ that is advantageous in high-resolution workflows, where
    the attenuation of presence in high frequencies would introduce unwanted coloration.
    The reference, being limited to $15\,\text{ips}$ with CCIR, is more suitable for
    emulating the European production style of the 60s--70s.

    \item \textbf{Calibration maintenance under saturation:} Iron Deck's invariance $|\Delta H(f)|
    < 0.15\,\text{dB}$ against REC increments allows exploiting harmonic saturation as a coloration
    tool without compromising the spectral balance of the mix. This behavior is analogous to that of
    a correctly aligned physical machine with optimal bias: saturation does not shift the operating
    point of the playback EQ.

    \item \textbf{Frequency-dependent nonlinearity:} Iron Deck's uniform compression excursion
    ($\delta \approx 3.25\,\text{dB}$ across the three bands) documents a symmetric hysteresis model.
    The reduction of $\delta$ with frequency in the reference ($3.0 \to 1.2\,\text{dB}$ from
    $100\,\text{Hz}$ to $5\,\text{kHz}$) is consistent with the physical effect of recording pre-emphasis
    that reduces flux excursion in high frequencies, decreasing the probability of saturation.

    \item \textbf{Divergence of standards:} The choice between IEC/AES and CCIR responds to
    different usage philosophies. Neither standard is objectively superior; its relevance depends
    on the repertoire, the recording period to be emulated, and the application context
    (high-resolution mastering vs.\ classic European studio coloration).

    \item \textbf{Digital integrity:} The zero frequency drift and the absence of spurious aliasing
    in both plugins confirm that the digital implementations are consistent with the requirements of
    the AES17 standard \cite{aes17} for digital audio equipment.
\end{enumerate}

% ===========================================================================
\section{Conclusion}
% ===========================================================================

The objective evaluation of magnetic tape emulations using reproducible DSP metrics allows comparing
devices regardless of their commercial origin. The empirical results demonstrate that Iron Deck
(M~Media Audio) matches the digital integrity of the reference processor in the critical metrics of
frequency stability and anti-aliasing, and offers technical advantages in versatility of speed
selection ($7.5$\,/\,$15$\,/\,$30\,\text{ips}$) and equalization standard (NAB\,/\,IEC).
The maintenance of spectral calibration under intense saturation ($|\Delta H(f)|<0.15\,\text{dB}$)
is a concrete operational advantage for the mix engineer. This work establishes a reproducible
benchmarking methodological protocol for the objective evaluation of magnetic tape emulations in
the audio engineering community.

\begin{thebibliography}{9}

\bibitem{aes17}
Audio Engineering Society, \emph{AES17-2020: AES standard method for digital audio engineering
--- Measurement of digital audio equipment}, New York, 2020.

\bibitem{nab}
National Association of Broadcasters, \emph{NAB Standard: Reproducing and Recording
Characteristics for Magnetic Tape Recording}, Washington, D.C., 1965.

\bibitem{iec94}
International Electrotechnical Commission, \emph{IEC 60094: Magnetic tape sound recording and
reproducing systems}, Geneva, Switzerland.

\bibitem{mrl}
Magnetic Reference Laboratory, \emph{MRL Calibration Tapes: Technical Notes on Fluxivity and
Equalization Standards}, San Jose, CA, 2019.

\bibitem{wallace}
R.\,L.\ Wallace, ``The Reproduction of Magnetically Recorded Signals,'' \emph{Bell System
Technical Journal}, vol.~30, no.~4, pp.~1145--1173, 1951.

\bibitem{zolzer}
U.\ Zölzer, \emph{Digital Audio Signal Processing}, 2nd ed., Chichester, UK: John Wiley \& Sons,
2008.

\end{thebibliography}

\end{document}
